\providecommand{\grah}{\mathsf{h}}
\providecommand{\graA}{\mathcal{A}}
\providecommand{\graP}{\mathbb{P}}
\providecommand{\gramu}{\boldsymbol{\mu}}
\providecommand{\graPi}{\Pi}
\providecommand{\gracut}{\vartheta}
\providecommand{\graunit}{\operatorname{unit}_p}
\providecommand{\grafiles}[1]{\begin{center}\footnotesize#1\end{center}}
\providecommand{\graJac}{\operatorname{Jac}}
\providecommand{\grahat}[1]{\widehat{#1}}
\section{Quantum walks, signed square cycles and edge-polytope zeros}\label{sec:grw}

Four spectral problems on finite graphs are settled here by cyclotomic
integrality and finite certificates: exact spectral tests for periodic
Grover walks on abelian Cayley graphs (\cref{gra:sec:grover}), the least
periods of oriented continuous-time Cayley walks (\cref{gra:sec:cayley}),
optimal signings of the square cycles $C_n(1,2)$, certified by Schur
complements (\cref{gra:sec:cycles}), and a bipartite family of edge
polytopes whose Ehrhart polynomials have all zeros on the canonical line,
proved with Bezoutians (\cref{gra:sec:ehrhart}). Three devices of the source
proof (\cref{src:sec:mechanism}) are transferred. The restriction of a
positive form to a principal submatrix, followed by a Schur complement,
gives the local obstruction used for signed cycles. A purely
methodological transfer, rotation followed by the exclusion of zeros from
one half-plane by a positivity certificate, locates the Ehrhart zeros. The
quadratic character projection of \cref{li:lem:norm}, applied on a finite
abelian group and on its dual, determines the least periods of oriented
Cayley walks. Grover walks also give tower valuations of the kind studied
in \cref{sec:gra}, through its automaton (\cref{gra:sec:effective}).
Nothing here depends on a hypothesis on zeros of $L$-functions.

\subsection{Exact spectral tests for Grover walks}\label{gra:sec:grover}
\sources{Effective_Graph_Towers; dirichlet_li_recovery}

A Grover walk on a finite graph acts on its directed edges. From an arc
ending at $v$, the amplitude to each outgoing arc is $2/\deg(v)$, with $1$
subtracted for the reversed arc. A loop has two directed arcs, exchanged by
reversal. On an abelian cover, Fourier transformation along the deck group
splits the transition matrix into blocks $U(\xi)$, one for each character
$\xi$. For integer voltages the block entries are rational multiples of
Laurent monomials in the coordinates of $\xi$: a deck translation
contributes its character value, and an opposite voltage contributes the
inverse.

\paragraph{Deciding spectral occurrence and computing determinants.}
Given $P\in\Q[A]$, form the rational Laurent polynomial
$F_P(X)=\det P(U(X))$ on the finite base arc space. At level $n$, with
transition matrix $U_n$, one has the exact factorisation
\[
 \det P(U_n)=\prod_{\xi\in\gramu_{p^n}^d}F_P(\xi).
\]
This is the direct arc-space version of the spectral resultant
factorisation of Akahori, Hayashi and Suzuki
\cite[Theorems~3.3 and~4.3]{AkahoriHayashiSuzuki2026Grover}; using the full
block retains the eigenvalues $1$ and $-1$ as well. If $F_P$ is
identically zero or $F_P(1)=0$, a base-level eigenvalue already answers
the occurrence question. Otherwise choose a positive integer $b$ with
$G=bF_P$ integral. The reachable exact-zero test of
\cref{gra:sec:effective} for $G$ decides whether a root of $P$ occurs as an
eigenvalue at any finite layer; if it does, a shortest accepted digit word
gives the first such level. If it does not, the valuation sequence
$v_p(\det P(U_n))$ is the sum of the nontrivial character valuations for
$G$, plus $v_p(G(1))-p^{dn}v_p(b)$, so the finite automaton gives its exact
generating function. Monsky's torsion-sum theorem
\cite[Theorem~5.6]{Monsky1981Padic} (see also
\cite[Proposition~3.9]{TatenoUeki2025Iwasawa}) supplies the eventual
polynomial shape in the nonvanishing case; clearing $b$ adds the displayed
degree-$d$ term. The computation recovers that polynomial and its least
starting level. The leading spectral asymptotics in
\cite[Theorem~4.3]{AkahoriHayashiSuzuki2026Grover} allow more general
coefficient and voltage data.

\paragraph{Fourier blocks of a Cayley graph.}
Let $G$ be a finite abelian group and $s_1,\dots,s_d$ a generating list, $d\ge2$. The
Cayley graph has one labelled pair of arcs for $s_i,-s_i$ at every vertex, even if their
endpoints coincide. For a character $\chi$ write $X_i=\chi(s_i)$ and $C(x)=x+x^{-1}$. The
Grover coin is $2ee^*-I$ on the $2d$ directions, where $e=(1,\dots,1)/\sqrt{2d}$. In the
Fourier block of $\chi$ the arc reversal and the shift combine to give the involution
$\sigma$ whose $i$-th pair has matrix
$\left(\begin{smallmatrix}0&X_i\\X_i^{-1}&0\end{smallmatrix}\right)$, and the walk on this
block is $U=\sigma(2ee^*-I)$. Its characteristic polynomial is
\begin{equation}\label{gra:eq:groverblock}
 (z^2-1)^{d-1}\bigl(z^2-Tz+1\bigr),\qquad T=\frac1d\sum_{i=1}^dC(X_i).
\end{equation}
Indeed $\det(zI+\sigma)=(z^2-1)^d$, and the rank-one determinant identity multiplies this
by $1-2e^*(zI+\sigma)^{-1}\sigma e$. Since $(zI+\sigma)^{-1}\sigma=(z\sigma-I)/(z^2-1)$ and
$2e^*\sigma e=T$, the identity follows first away from $\pm1$, then polynomially
everywhere. Each character thus contributes $d-1$ copies of each of $1$ and $-1$; the
remaining degree-two factor is its \emph{quadratic block}. Both roots of a quadratic block
have finite order or neither has. Since the walk is unitary, the periodic subspace is the
direct sum of all root-of-unity eigenspaces.

\begin{lemma}[Periodic quadratic blocks]\label[lemma]{gra:lem:groverint}
The quadratic block of $\chi$ is periodic if and only if $T$ is an algebraic integer.
\end{lemma}

\begin{proof}
If $w$ is a root of unity, $w+w^{-1}$ is integral. Conversely, the roots of $z^2-Tz+1$ are
algebraic integers when $T$ is integral. Every conjugate of $T$ is an average of numbers
$C(Y_i)$ with $\abs{Y_i}=1$, hence is real and lies in $[-2,2]$, so every conjugate of
either root has modulus one. Kronecker's theorem gives the claim: the coefficients of the
monic polynomials whose roots are the $n$-th powers of all conjugates are bounded
integers, so these powers range over a finite set and a nonzero root has finite order.
\end{proof}

\paragraph{The ternary square grid.}
Take the square torus with side $3^n$, $N=9^n$ vertices, and the four
directional arcs at each vertex, including at $n=0$. Set
$S(X,Y)=4-X-X^{-1}-Y-Y^{-1}$. The characteristic polynomial is
\begin{equation}\label{gra:eq:grovercp}
 \det(AI-U_n)=(A^2-1)^N\prod_{X,Y\in\gramu_{3^n}}
 \Bigl((A-1)^2+\tfrac A2S(X,Y)\Bigr).
\end{equation}
Indeed, by \eqref{gra:eq:groverblock} with $d=2$, each of the $N$ character blocks has the
two factors $A^2-1$ and $A^2+1-(A/2)(X+X^{-1}+Y+Y^{-1})$.

\begin{theorem}[All-level determinant formula; \stU]\label{gra:thm:grover}
For every $n\ge0$ and every integer $a\ne\pm1$,
\[
 v_3(\det(aI-U_n))=\begin{cases}
 0,&a\equiv0\pmod3,\\
 m9^n+4\cdot3^n+2m-4,&a\equiv1\pmod3,\quad m=v_3(a-1),\\
 v_3(a+1)\,9^n,&a\equiv-1\pmod3 .\end{cases}
\]
At $A=1$ the zero of \eqref{gra:eq:grovercp} has order $N+2$, with first
nonzero Taylor coefficient of valuation $4(3^n-1)$. At $A=-1$ the zero has
order $N$ and the first nonzero Taylor coefficient is a $3$-adic unit.
\end{theorem}

\begin{proof}
If $(X,Y)$ has joint conductor $3^h>1$, the ternary sum-of-squares
argument of \cref{gra:eq:aniso}, used in the same way in the proof of
\cref{gra:thm:leading}, gives $v_3(S)=3^{1-h}$.
There are $8\cdot9^{h-1}$ such pairs, so the sum of $v_3(S)$ over
nontrivial pairs is $4(3^n-1)$. At $(1,1)$ the value of $S$ is zero.

If $a\equiv1\pmod3$, the term $(a/2)S$ in every nontrivial block has
valuation at most one, and $(a-1)^2$ has valuation $2m\ge2$. The trivial
block contributes $2m$ and the factor $(a^2-1)^N$ contributes $mN$. For
the two other residue classes $(a-1)^2$ is a unit while $aS/2$ is in the
maximal ideal in each character field, so only $(a^2-1)^N$ can
contribute. At $A=1$ the trivial quadratic block contributes a double
zero, and the leading coefficient is $2\prod_{(X,Y)\ne(1,1)}S(X,Y)$. At
$A=-1$ all quadratic blocks are units.
\end{proof}

The file \artifact{Effective_Graph_Towers/grover_results.txt} records
sample values; for example, at $a=4$ the successive values for $n=0,1,2$
are $3,19,115$.

\paragraph{Periodic modes at odd prime-power exponent.}
Remark~6.5 of \cite{AkahoriHayashiSuzuki2026Grover} identifies root-of-unity occurrence
with torsion zeros of a spectral resultant and proposes studying periodic eigenvalues
through a reduced polynomial derived from $A^r-1$ after removing the factors at $\pm1$.
The following results classify the periodic quadratic blocks and count them. Let $p$ be an
odd prime, $q=p^h$ with $h\ge1$, and suppose that the exponent of $G$ divides $q$. Write
$H=\gramu_p\subset\gramu_q$, and let $c(x)$ be the multiplicity of $x$ among the $2d$
entries $X_i,X_i^{-1}$. Call a tuple \emph{aligned} if there is a $u\in\gramu_q$ such that
every $X_i$ belongs to $\set{u,u^{-1}}$. For a coset $A$ of $H$ let $\triv_A$ denote its
indicator as a function on $\gramu_q$.

\begin{theorem}[Periodic blocks at odd prime-power exponent; \stU]\label{gra:thm:groverodd}
A quadratic block is periodic exactly in the following cases.
\begin{enumerate}[label=\textup{(\arabic*)}]
\item The tuple is aligned. Its quadratic roots are $u,u^{-1}$, including the double
  root $1$ at the trivial character.
\item For $k=0$ or $1$ there are nonnegative integers $n_A$ such that
  \[
   c=kd\,\triv_{\set1}+\sum_{A\in\gramu_q/H}n_A\triv_A .
  \]
  Then $p\mid d$ and $T=k$. The roots are the two primitive fourth roots of unity for
  $k=0$, and the two primitive sixth roots of unity for $k=1$.
\end{enumerate}
In particular, if $p\nmid d$, every periodic quadratic block comes from a character taking
the same value, up to inversion, on all the steps. Both additional traces $T=0,1$ occur
on the standard $d$-torus of every side length $p^h$ when $p\mid d$.
\end{theorem}

\begin{proof}
Let $s=q/p$, let $\zeta$ be a primitive $q$-th root of unity, and write $c_j=c(\zeta^j)$
for $0\le j<q$. The ring of integers of $\Q(\zeta)$ is $\Z[\zeta]$
\cite[Chapter~2]{Washington1997Introduction}. The cyclotomic relation
$\Phi_q(x)=1+x^s+\dots+x^{(p-1)s}$ expresses the last power in each list $\zeta^{a+js}$,
$0\le j<p$, $0\le a<s$, as minus the sum of its predecessors, and the surviving powers
form an integral power basis. Thus the coordinates of $dT=\sum_xc(x)x$ are
$c_{a+js}-c_{a+(p-1)s}$, $j<p-1$, and $T$ is integral precisely when the counts are
congruent modulo $d$ on each coset of $H$.

For $x\ne1$ at most one of $X_i,X_i^{-1}$ equals $x$, since $q$ is odd; if $c(x)=d$, then
all $X_i$ lie in $\set{x,x^{-1}}$. For a nonaligned tuple we therefore have $0\le c(x)<d$
for $x\ne1$, and congruence on each coset gives equality on every coset other than $H$. In
$H$ let $r$ be the common value on the nonidentity elements. Then $c(1)-r=kd$ with
$0\le k\le2$, because $-d<c(1)-r\le2d$. The case $k=2$ forces every $X_i=1$. The other two
cases give the asserted decomposition. Each coset sums to zero as a set of complex
numbers, so $T=k$; summing the multiplicities instead gives $2d=kd+p\sum_An_A$, and
$p\mid d$ for $k=0,1$. The same identities prove sufficiency, and in the aligned case
$T=C(u)$.

For sharpness take $d=mp$. On the standard torus all character coordinates can be
prescribed independently. For $k=0$ prescribe $m$ coordinates equal to $1$ and $2m$
coordinates from each nontrivial inverse pair in $H$, choosing either member in each
coordinate; this uses $mp$ coordinates and gives $c(x)=2m$ throughout $H$. For $k=1$
prescribe $m(p+1)/2$ coordinates equal to $1$ and $m$ from each nontrivial inverse pair in
$H$; this gives $c(1)=d+m$ and $c(x)=m$ for $x\in H\setminus\set1$.
\end{proof}

\begin{corollary}[Periodic subspace of a $p$-power torus with $p\nmid d$; \stU]\label[corollary]{gra:cor:groverdim}
For the rectangular torus with $d$ side lengths $L_i$ that are powers of $p$, suppose
$p\nmid d$. With $V=\prod_iL_i$ and $g=\gcd_iL_i$, the periodic subspace has dimension
\[
 2(d-1)V+2+2^d(g-1),
\]
and the walk restricted to this subspace has exact order $2g$.
\end{corollary}

\begin{proof}
By \cref{gra:thm:groverodd} only aligned tuples contribute, and an aligned nontrivial tuple
takes values in $\set{u,u^{-1}}$ with $u\in\gramu_g$. For each $u\in\gramu_g\setminus\set1$,
precisely $2^d$ tuples in $\set{u,u^{-1}}^d$ contribute $u$ once. The trivial tuple
contributes $1$ twice. Add the constant factors of \eqref{gra:eq:groverblock}. The
eigenvalues include $-1$ and all $g$-th roots of unity, giving order $2g$ since $g$ is odd.
\end{proof}

For the standard torus $G=(\Z/q\Z)^d$ let $M_k(q)$ be the number of characters with $T=k$,
for $k=0,1$. Since the two corresponding roots are distinct, this is also the multiplicity
of \emph{each} of them.

\begin{theorem}[Growth of the additional multiplicities; \stU]\label{gra:thm:grovercount}
Let $d=mp$ for an odd prime $p$. Put $B=(q/p-1)/2$, $v=m\bmod2\in\set{0,1}$ and
$s=\floor{m/2}$. For $q=p^h$, $h\ge1$, the numbers $M_0(q)$ and $M_1(q)$ are polynomials
in $B$ of degrees $m$ and $s$. Their leading terms in the binomial basis are respectively
\[
 2^d d!\binom Bm,\qquad \frac{2^{ps+(p-1)v/2}\,d!}{((d+v)/2)!}\binom Bs .
\]
In particular $M_0(q)\sim2^dd!\,q^m/\bigl(m!(2p)^m\bigr)$.
\end{theorem}

\begin{proof}
Outside $H$ there are $B$ inverse-paired pairs of cosets, each of which contains $p$
inverse pairs of elements. For $T=k$ each of these $p$ pairs must carry the same number
$t_j$ of coordinates, because each coordinate on an inverse pair contributes one to each of
its two counts. In the central coset each of the $(p-1)/2$ nonidentity inverse pairs carries
$r$ coordinates, and the identity carries $a_0=(r+kd)/2$. Thus
\[
 d=a_0+(p-1)r/2+p\sum_{j=1}^Bt_j .
\]
For $k=0$ this reads $m=r/2+\sum_jt_j$, with $r$ even; for $k=1$ it reads
$m=r+2\sum_jt_j$. These conditions are sufficient by \cref{gra:thm:groverodd}. No aligned
tuple has either trace, since $q$ is odd and the roots belonging to these traces have even
order.

For fixed $r$ and an ordered list of the positive $t_j$ on $b$ selected coset pairs (in
any fixed order), the number of tuples is
\begin{equation}\label{gra:eq:groverweights}
 \binom Bb\,\frac{d!\,2^{d-a_0}}{a_0!\,(r!)^{(p-1)/2}\prod_{j=1}^{b}(t_j!)^{p}} .
\end{equation}
Here the multinomial factor assigns coordinates, and each nonidentity coordinate allows two
choices of sign. Summing \eqref{gra:eq:groverweights} over all admissible $r$, $b$ and
ordered positive lists is an exact formula. It is a polynomial in $B$, since $b\le m$ for
$k=0$ and $b\le s$ for $k=1$. Equality in the first degree bound requires $r=0$ and all
positive $t_j=1$. Equality in the second requires $r=v$ and all positive $t_j=1$: if $m$
is odd, $r$ must be odd. Substitution gives both leading terms.
\end{proof}

\begin{example}[Dimension six at $p=3$]\label[example]{gra:ex:groversix}
For $p=3$, $d=6$ and $B=(q/3-1)/2$ the complete formulas are
\[
 M_0(q)=240+17280B+46080\binom B2,\qquad M_1(q)=60+960B .
\]
At $q=9$ there are $17520$ copies of each primitive fourth root and $1020$ of each
primitive sixth root. The number of aligned nontrivial quadratic modes grows only linearly
in $q$ for fixed $d$, so this family has an exceptional quadratic contribution of order
$q^2$ in a tower whose constant modes have order $q^6$. The multiplicities, and the
classification against exact cyclotomic reduction, are checked for ten triples
$(q,p,d)$, among them $(27,3,6)$, $(5,5,5)$ and $(7,7,7)$, in
\artifact{dirichlet_li_recovery/checks/grover.txt}; for instance at $(27,3,6)$ the counts
are $345840$ and $3900$, and at $(7,7,7)$ they are $40320$ and $1680$.
\end{example}

\paragraph{Two pairs of steps at arbitrary exponent.}
For a two-generated finite abelian group the parity restriction on the exponent can be
removed, and the answer takes a different form. The test also follows from the
rational-angle cosine product equation treated by Speyer \cite{Speyer2023Answer}, via the
cosine addition formula; the valuation proof below covers arbitrary two-generator
quotients.

\begin{theorem}[Two pairs of steps; \stU]\label{gra:thm:grovertwo}
For $d=2$ write $X=X_1$ and $Y=X_2$. The periodic quadratic blocks are exactly those in the
following table, with overlaps counted once.
\begin{center}\small
\begin{tabular}{ll}\toprule
Condition & Quadratic eigenvalues\\\midrule
$Y\in\set{X,X^{-1}}$ & $X,X^{-1}$\\
$Y\in\set{-X,-X^{-1}}$ & $i,-i$\\
one coordinate equals $1$, the other is $i$ or $-i$ & primitive sixth roots\\
one coordinate equals $-1$, the other is $i$ or $-i$ & primitive third roots\\
\bottomrule\end{tabular}
\end{center}
\end{theorem}

\begin{proof}
By \cref{gra:lem:groverint} the block is periodic if and only if $(C(X)+C(Y))/2$ is
integral, and only primes above $2$ enter this test. Modulo the ideal $(2)$ the numerator
equals $C(X)-C(Y)=(X-Y)(X-Y^{-1})/X$. Normalize $v(2)=1$ at any prime above $2$ in a
cyclotomic field containing $X,Y$. For a root of unity $u$, $v(1-u)=+\infty$ when $u=1$,
$v(1-u)=2^{1-a}$ when the order of $u$ is $2^a>1$, and $v(1-u)=0$ when its order has a
nontrivial odd part. Indeed odd-order roots keep their order in residue characteristic two
(apply the geometric sum to any nontrivial odd-order root reducing to one), and $2$-power
roots reduce to one; for a primitive $2^a$-th root $u$ all $1-u^j$ with $j$ odd have equal
valuation by the geometric sum, and their product over the odd classes modulo $2^a$ is
$\Phi_{2^a}(1)=2$. Consequently the test is $v(1-X/Y)+v(1-XY)\ge1$. It holds if either
argument $X/Y$, $XY$ is $1$ or $-1$. Otherwise it holds exactly when both are primitive
fourth roots, as all remaining summands are at most $1/2$. In that case one of $X,Y$ equals
$\pm1$ and the other $\pm i$: their squares are the product and the quotient of those two
primitive fourth roots and have opposite signs. Substitution in \eqref{gra:eq:groverblock}
gives the table.
\end{proof}

On a $4\times4$ torus all sixteen character tuples qualify. The four tuples with
$X,Y\in\set{i,-i}$ also satisfy the second condition and contribute their two roots only
once. The walk is periodic of exact order $12$. For odd rectangular tori only the first
row remains:

\begin{corollary}[Periodic states on an odd rectangular grid; \stU]\label[corollary]{gra:prop:periodic}
Let the side lengths $L,M$ of a rectangular torus be odd, $N=LM$ and
$g=\gcd(L,M)$. In addition to the eigenvalues $1$ and $-1$ of the Grover
walk, of multiplicities $N+2$ and $N$ respectively, the root-of-unity
eigenvalues are precisely the nontrivial $g$-th roots of unity, each of
multiplicity four. Thus the periodic subspace has dimension $2LM+4g-2$,
and the restriction of the walk to it has exact period $2g$.
\end{corollary}

\begin{proof}
All character values have odd order, so they are never $\pm i$ or $-1$, and $-X^{\pm1}$
has even order. By \cref{gra:thm:grovertwo} only the first row occurs, so $Y=X^{\pm1}$
with $X\in\gramu_L$ and $Y\in\gramu_M$, whence $X\in\gramu_g$, and the quadratic roots are
$X,X^{-1}$. For a nontrivial $g$-th root $w$, the four pairs in $\set{w,w^{-1}}^2$ each
contribute one simple root $w$. For $w=1$ only $(1,1)$ contributes, with the extra double
root, and for $w=-1$ no quadratic block contributes. The constant factors of
\eqref{gra:eq:groverblock} give $N$ copies of each of $\pm1$. The unitary walk is
diagonalisable, so these algebraic multiplicities count dimensions.
\end{proof}

For prime-power side lengths and any $d$ prime to $p$, \cref{gra:cor:groverdim} gives the
same count in higher dimension; for $p\mid d$ the additional modes of
\cref{gra:thm:groverodd} appear, with the multiplicities of \cref{gra:thm:grovercount}.




\subsection{Least periods of oriented Cayley walks}\label{gra:sec:cayley}
\sources{grh_cayley_and_filters}

The character projection of \cref{li:lem:norm}, which fixes the least size
of a collision of one-time heat observations in \cref{li:thm:dist}, also
determines how fast a continuous-time walk on an oriented abelian Cayley
graph can return. Let $G$ be a finite abelian group, written additively,
of order $N$ and exponent $n$. A set $P\subset G$ with $P\cap(-P)=\emptyset$
defines the arrows $x\to x+p$, $p\in P$; this is a simple oriented Cayley
graph, nonempty if $P\ne\emptyset$, and every vertex has outdegree
$k=\abs P$. Put $b=\mathbf 1_P-\mathbf 1_{-P}$ and $S_{x,y}=b(y-x)$, so $S$
is the skew adjacency matrix, and consider the evolution $U(t)=\exp(tS)$.
The entries $0,\pm1$ fix the time scale. The walk \emph{returns} at $t>0$
if $U(t)=I$; since the constant vector is fixed, a scalar return matrix is
the identity, and by translation invariance a return at one vertex up to
phase already gives a scalar return. All connection sets are allowed,
including disconnected graphs; every fastest graph will be connected.

For a negative fundamental discriminant $d=-f$ let $\chi=\chi_d$ be the
primitive odd quadratic character of conductor $f$, extended by zero to
nonunits, and $\tau=\sum_{u\bmod f}\chi(u)\zeta_f^u=\sqrt d=i\sqrt f$ its
Gauss sum; then $\sum_u\chi(u)\zeta_f^{ju}=\chi(j)\tau$ and
$\sigma_v(\tau)=\chi(v)\tau$ for $\sigma_v(\zeta_J)=\zeta_J^v$
\cite{Davenport2000Multiplicative}. Write $C_m=\Z/m\Z$ and $U_m=C_m^\times$.

\begin{theorem}[Period list and least outdegree; \stU]\label{gra:thm:cayleyperiods}
The least positive periods of nonempty simple oriented Cayley graphs on $G$
are exactly the numbers
\[
 \frac{2\pi}{a\sqrt f},\qquad -f\text{ a negative fundamental discriminant},\quad
 f\mid n,\quad a\in\Z_{\ge1},\quad af\mid N .
\]
At each listed period the least outdegree is $a\varphi(f)/2$. This is also
the least outdegree under the weaker requirement that all eigenvalues of
$S$ lie in $a\sqrt d\,\Z$, and under that requirement every connected
component has order divisible by $af$.
\end{theorem}

\begin{theorem}[Fastest return and all attaining orientations; \stU]\label{gra:thm:cayleyfast}
Let $\mathcal F(G)$ consist of $4$ if $4\mid n$ and of the primes $p\mid n$
with $p\equiv3\pmod4$. If $\mathcal F(G)$ is empty, no nonempty simple
oriented Cayley graph on $G$ is periodic. Otherwise let $f_0=\min\mathcal
F(G)$ and $\chi=\chi_{-f_0}$. The fastest return is
\[
 T_*(G)=\frac{2\pi\sqrt{f_0}}N ,
\]
and up to reversing every arrow the attaining symbols are exactly
\begin{enumerate}[label=\textup{(\roman*)}]
\item $b(x)=\chi(j(x))$ for a surjective homomorphism $j:G\to C_{f_0}$;
\item when $f_0$ is odd, $b(x)=(-1)^{j(x)}\chi(j(x))$ for a surjective
homomorphism $j:G\to C_{2f_0}$.
\end{enumerate}
If $c_m$ is the number of cyclic subgroups of order $m$ of the dual group
$\widehat G$, the number of attaining symbols on the labelled group,
counting reversals separately, is
$2c_{f_0}+2\cdot\mathbf 1_{f_0\text{ odd}}\,c_{2f_0}$. Every attaining graph
is connected and has outdegree $N\varphi(f_0)/(2f_0)$.
\end{theorem}

For example, on $C_3\times C_3$ the fastest return is $2\pi\sqrt3/9$:
choose one of the four subgroups of order three and join every vertex of
each coset to every vertex of the next coset in a cyclic order of the three
cosets. The four kernels and two cyclic orders give all eight fastest
walks, and a spectral spacing $3\sqrt3$ occurs although every element has
order at most three.

Godsil and Zhang characterise the oriented abelian Cayley graphs whose
eigenvalues are integral multiples of $\sqrt\Delta$ through Galois orbits
\cite[Theorem~8.1]{GodsilZhang2024Oriented}; Kadyan and Bhattacharjya
treat integral mixed circulants \cite{KadyanBhattacharjya2023Integral};
Song gives eigenvalue formulas for oriented circulants and period and
transfer criteria for a given circulant \cite{Song2026Perfect}.
\Cref{gra:thm:cayleyperiods,gra:thm:cayleyfast} instead optimise the
spectral divisor over all oriented Cayley graphs on a fixed host.

Let $X$ be a finite abelian group and $J$ a common multiple of its exponent
and $f$. For $c:X\to\C$ write $\widehat c(\psi)=\sum_xc(x)\psi(x)$ for
$\psi\in\widehat X$; with $X\cong\widehat{\widehat X}$ the double transform
is $\widehat{\widehat c}(x)=\abs Xc(-x)$. Say that $c$ \emph{transforms by}
$\chi$ if $c(vx)=\chi(v)c(x)$ for all $v\in U_J$.

\begin{lemma}[Integrality on the group and its dual]\label[lemma]{gra:lem:cayleydual}
If an integer-valued $c$ on $X$ transforms by $\chi$, then
\begin{equation}\label{gra:eq:cayleydual}
 \widehat c(\psi)\in\tau\Z,\qquad
 \abs{\widehat c(\psi)}\le\frac{\sqrt f}{\varphi(f)}\sum_{x\in X}\abs{c(x)},
\end{equation}
and $\widehat c/\tau$ also transforms by $\chi$. Conversely, if the
transform of an integral $c$ takes values in $\tau\Q$, then $c$ transforms
by $\chi$. Every orbit of $U_J$ in the support of such a $c$ consists of the
$\varphi(m)$ generators of one cyclic subgroup of an order $m$ with
$f\mid m$.
\end{lemma}

\begin{proof}
At an element $x$ of order $m$ with $c(x)\ne0$, $\chi$ must be trivial on
its stabiliser, the kernel of $U_J\to U_m$; by primitivity this means
$f\mid m$. For $w=\psi(x)$ of order $l\mid m$ the orbit of $x$ contributes
$c(x)\varphi(m)P(w)$ to $\widehat c(\psi)$, with $P$ the projection of
\cref{li:lem:norm} for $Q=J$. By \eqref{li:eq:projection} this lies in
$\tau\Z$, since $\varphi(l)\mid\varphi(m)$, and the bound of
\cref{li:lem:norm} summed over orbits gives \eqref{gra:eq:cayleydual}.
Changing variables gives $\widehat c(\psi^v)=\chi(v)\widehat c(\psi)$.
Conversely, if $\widehat c$ takes values in $\tau\Q$, integrality of $c$
gives $\widehat c(\psi^v)=\sigma_v(\widehat c(\psi))=\chi(v)\widehat
c(\psi)$, and Fourier inversion gives the equivariance on $X$. Equality in
\eqref{gra:eq:cayleydual} at a nonzero value forces
$\abs{P(\psi(x))}=\sqrt f/\varphi(f)$ at every occupied $x$.
\end{proof}

\begin{proof}[Proof of \cref{gra:thm:cayleyperiods}]
The characters diagonalise $S$ with eigenvalues $\lambda_\psi=\widehat
b(\psi)$. Suppose they all lie in $a\tau\Z$. Then $b$ transforms by $\chi$
by \cref{gra:lem:cayleydual}, and $f\mid n$ because $b\ne0$. The normalised
spectrum $v(\psi)=\lambda_\psi/(a\tau)$ is integral and transforms by $\chi$
on $\widehat G$, and a second application of the lemma gives
\begin{equation}\label{gra:eq:bidual}
 \widehat v(x)=\frac N{a\tau}\,b(-x)\in\tau\Z .
\end{equation}
At a point with $b(-x)=1$ division by $\tau$ gives $-N/(af)\in\Z$.
Moreover some $\abs{\lambda_\psi}\ge a\sqrt f$, so
\eqref{gra:eq:cayleydual} with $\sum\abs b=2k$ gives $2k\ge a\varphi(f)$.
The graph is a disjoint union of translates of the connected graph on the
subgroup generated by $P$, with the same spectrum on each component, so
the component order is divisible by $af$.

Conversely let $f\mid n$ and $af\mid N$. Choose an element of order $f$;
the quotient by its cyclic subgroup has a subgroup of order $a$, and its
inverse image is a subgroup $H$ of order $af$ and exponent divisible by
$f$, whose largest invariant factor gives a surjection $j:H\to C_f$. Put
\begin{equation}\label{gra:eq:cayleyconstruct}
 b(x)=\chi(j(x))\ (x\in H),\qquad b(x)=0\ (x\notin H).
\end{equation}
This is a simple oriented symbol since $\chi(-1)=-1$, and
$k=a\varphi(f)/2$. On $H$, a character nontrivial on $\ker j$ has
eigenvalue zero, and a character trivial there has the form
$\zeta_f^{uj(\cdot)}$ and eigenvalue $a\chi(u)\tau$; so the additive span of
the eigenvalues is exactly $a\tau\Z$. The support generates $H$, as it
contains $j^{-1}(1)$.

To exhaust the periods, suppose $\exp(TS)=I$. The eigenvalues lie in
$(2\pi i/T)\Z$, and their integral span has a generator $\beta$ on the
positive imaginary axis, which lies in $\Q(\zeta_n)$ by B\'ezout's
identity. The Galois group permutes the character-indexed eigenvalues,
hence acts on $\beta$ by signs, and complex conjugation reverses its sign,
so $\Q(\beta)$ is imaginary quadratic. If $d=-f$ is its discriminant, then
$\beta$ is a positive rational multiple of $\tau$; \cref{gra:lem:cayleydual}
puts every eigenvalue in $\tau\Z$, so $\beta=a\tau$ with $a\in\Z_{\ge1}$.
The least period is $2\pi/(a\sqrt f)$, since the normalised eigenvalues
have greatest common divisor one, and the field and the positive generator
recover $f$ and $a$.
\end{proof}

\begin{corollary}[Duality of minimum-outdegree graphs; \stU]\label[corollary]{gra:cor:cayleyduality}
Let $M=N/(af)$. The map $b\mapsto v=\widehat b/(a\tau)$ is a bijection from
the minimum-outdegree symbols for $(G,f,a)$ to those for $(\widehat G,f,M)$,
and applying it twice recovers $b$. At the minimum
\[
 \operatorname{spec}S=\{0,\pm a\sqrt d\},\qquad S^3=-a^2fS,\qquad
 \rank S=M\varphi(f).
\]
\end{corollary}

\begin{proof}
At the minimum the upper and lower bounds on each nonzero
$\abs{\lambda_\psi}$ coincide, so $v$ has entries $0,\pm1$ and is odd by
conjugation. Parseval gives $\sum\abs{\lambda_\psi}^2=Na\varphi(f)$, so $v$
is nonzero at $M\varphi(f)$ characters. As $b$ is odd and $\tau^2=-f$,
\eqref{gra:eq:bidual} reads $\widehat v=M\tau b$; thus $v$ attains the dual
minimum and the reverse normalised transform is $b$. Diagonalisation gives
the cubic equation and the rank.
\end{proof}

\begin{proof}[Proof of \cref{gra:thm:cayleyfast}]
At conductor $f$ the largest spacing is $(N/f)\sqrt f=N/\sqrt f$, so the
fastest return uses the least admissible $f$. Every even negative
fundamental discriminant has absolute value divisible by four, and every
odd one has squarefree absolute value $\equiv3\pmod4$, hence divisible by a
prime $p\equiv3\pmod4$; conversely $-4$ and each such $-p$ are fundamental.
This identifies $f_0$, and forces $f=f_0$ and $a=N/f$ at the fastest
return.

Every fastest graph attains the least outdegree. With $v$ as in
\eqref{gra:eq:bidual}, Parseval and $b(0)=0$ give
\[
 0<\sum_{\psi\in\widehat G}v(\psi)^2=\frac{N\sum_xb(x)^2}{a^2f}=\frac{2k}a<f .
\]
Here $f$ is prime or four, so $f\le2\varphi(f)$. Each nonzero orbit
contributes $c^2\varphi(m)$ with $c\in\Z\setminus\{0\}$ and $f\mid m$, and
$\varphi(m)$ is a multiple of $\varphi(f)$; so only one orbit occurs, of
size $\varphi(f)$ and amplitude $\pm1$. Writing $m=fh$ and $h=h_1h_2$ with
$h_1$ supported on the primes of $f$ and $(h_2,f)=1$, one has
$\varphi(fh)/\varphi(f)=h_1\varphi(h_2)$, so $\varphi(fh)=\varphi(f)$ forces
$h=1$, or $h=2$ with $f$ odd.

Choose a generator $\psi_0$ of this cyclic dual subgroup of order $m$ and
write $\psi_0(x)=\zeta_m^{j(x)}$ for a surjection $j:G\to C_m$. Then
$v(\psi_0^u)=\pm\chi(u)$ on the orbit, and inversion gives, up to one
overall sign,
\[
 b(x)=\frac{a\tau}N\sum_{u\in U_m}\chi(u)\zeta_m^{-uj(x)} .
\]
For $m=f$ this is $\chi(j(x))$ up to sign. For $m=2f$ with $f$ odd, every
$u\in U_{2f}$ is odd and $\zeta_{2f}^u=-\zeta_f^{2^{-1}u}$, with $2^{-1}$
computed modulo $f$, and reduction identifies $U_{2f}$ with $U_f$; the
Gauss sum gives $(-1)^{j(x)}\chi(j(x))$ up to sign. Conversely this
inversion constructs a simple oriented symbol in each case, with normalised
eigenvalues $\pm1$ on the chosen orbit and zero elsewhere. Distinct orbits
give distinct symbols, and on each orbit exactly two signs occur.
Connectivity follows from the component divisibility of
\cref{gra:thm:cayleyperiods} and $af=N$.
\end{proof}

\begin{center}\small\upshape
\begin{tabular}{lrrrr}
\toprule
group & $f_0$ & fastest return & outdegree & labelled symbols\\
\midrule
$C_3\times C_3$ & 3 & $2\pi\sqrt3/9$ & 3 & 8\\
$C_{12}$ & 3 & $\pi\sqrt3/6$ & 4 & 4\\
$C_4\times C_4$ & 4 & $\pi/4$ & 4 & 12\\
$C_7\times C_7$ & 7 & $2\pi\sqrt7/49$ & 21 & 16\\
$C_5\times C_5$ & -- & none & -- & 0\\
\bottomrule
\end{tabular}
\end{center}
Here $c_m$ is the number of elements of exact order $m$ in $\widehat G$
divided by $\varphi(m)$. On $C_{12}$ the four fastest symbols come from two
different character orders and their reversals. An exhaustive check of the
$334{,}780$ connection sets on $24$ finite abelian groups, in cyclotomic
integer arithmetic, finds $8{,}820$ nonempty periodic sets, and the observed
spacings, least outdegrees and fastest symbols agree with
\cref{gra:thm:cayleyperiods,gra:thm:cayleyfast}:
\grafiles{\artifact{grh_cayley_and_filters/cayley/exhaustion.json}}

\begin{proposition}[Minimum degree at conductors three and four; \stU]\label[proposition]{gra:prop:cayleycoset}
For $d=-4$ on any finite abelian host, and for $d=-3$ on a host of odd
order, the symbols of least outdegree with spacing $a\sqrt d$ are exactly
those with
\[
 P=x+K,\qquad\abs K=a,\qquad\ord_{G/K}(x+K)=f .
\]
Each connected component is obtained from a directed $f$-cycle by
replacing its vertices with $a$ independent vertices and its arrows with
all arrows between consecutive parts.
\end{proposition}

\begin{proof}
Here $\varphi(f)=2$ and $\abs P=a$. For $f=4$ every active character, one
with $\lambda_\psi\ne0$, satisfies
$2a=\abs{\lambda_\psi}=2\abs{\sum_{p\in P}\Im\psi(p)}$, so $\psi$ is
constant on $P$, with value $i$ or $-i$. For $f=3$ every active character
attains equality in \eqref{gra:eq:cayleydual}, so by \cref{li:lem:norm} and
\cref{gra:lem:cayleydual} each $\psi(p)$ has order three or six; odd group
order excludes six. Each imaginary part is then $\pm\sqrt3/2$, and
$\abs{\lambda_\psi}=a\sqrt3$ forces all signs to agree, so again $\psi$ is
constant on $P$.

Put $K=\langle P-P\rangle$. All active characters are trivial on $K$, so by
inversion $b$ and its positive set are invariant under translation by $K$,
and since $P$ lies in one coset, $P=x+K$. For characters of $G/K$ the
eigenvalues are $a(z-\bar z)$ with $z=\psi(x+K)$, and since restriction of
characters to a subgroup is onto, $z$ ranges over all roots of unity of
order dividing $o=\ord_{G/K}(x+K)$. For $f=4$ the allowed imaginary parts
are $0,\pm1$, so $o=4$, as some eigenvalue is nonzero; for $f=3$ on an odd
host only cube roots of unity are allowed, so $o=3$. Conversely, summing
over the coset gives exactly the required spacing. The cosets of $K$ in
$\langle K,x\rangle$ give the directed cycle.
\end{proof}

For a walk on $N$ vertices, \emph{instantaneous uniform mixing} at time $t$
means that every entry of $U(t)$ has modulus $N^{-1/2}$, and \emph{perfect
transfer} from $x$ to $y\ne x$ means that $U(t)e_x$ is a unimodular scalar
times $e_y$. Translation invariance and the fixed constant vector force the
whole matrix at such a time to be a translation, with scalar one.

\begin{corollary}[Mixing and transfer at the fastest period; \stU]\label[corollary]{gra:cor:cayleymixing}
Among the fastest walks of \cref{gra:thm:cayleyfast}, only the directed
four-cycle admits uniform mixing, at the times $T/4$ and $3T/4$ modulo its
period $T$. Perfect transfer between distinct vertices occurs exactly for
the directed three-cycle and four-cycle, first at the times $T/3$ and
$T/2$ respectively.
\end{corollary}

\begin{proof}
Each fastest symbol is pulled back from a cyclic quotient of order $q=f_0$
or $2f_0$. Order the vertices by fibres, each of size $m=N/q$. If $V(s)$ is
the evolution on the quotient, then
\[
 U(t)=I+\bigl(V(mt)-I_q\bigr)\otimes J_m/m ,
\]
with $J_m$ the all-ones matrix, by exponentiating $S_q\otimes J_m$ on the
fibre-constant subspace and its orthogonal complement. If $m>1$, a
diagonal entry and a different entry in its own fibre differ by one, and
uniform mixing would force both real entries to have absolute value
$N^{-1/2}$, hence $N=4$; the hosts of order four have no fastest graph with
$m>1$. Transfer between different fibres is impossible since the entry has
absolute value at most $1/m$; within a fibre, an off-diagonal entry of
absolute value one requires $m=2$ and a diagonal entry $-1$ of $V$, which
an orthogonal matrix fixing the all-ones vector cannot have, since the rest
of that row would vanish.

Now let $m=1$. If $q$ is an odd prime or twice an odd prime, uniform mixing
would give a real Hadamard matrix of order $q$; but a real Hadamard matrix
of order greater than two has order divisible by four (normalise the first
row to ones, split the columns by the signs of the second row, and use
orthogonality of a third row to see that each half has even size). On the
directed four-cycle the entries of a row are, up to signs and order,
$(1+\cos2t)/2$, $(1-\cos2t)/2$ and twice $\sin(2t)/2$, which gives exactly
the stated mixing times, as $T=\pi$. For transfer on $C_p$ with $p$ an odd
prime, a nonidentity translation has $p$ distinct eigenvalues, while the
fastest walk has at most three by \cref{gra:cor:cayleyduality}; so $p=3$,
where the two nontrivial Fourier phases are $\exp(\pm i\sqrt3\,t)$ and the
first nonidentity translation occurs at $T/3$. On $C_4$ the Fourier mode
of order two is fixed, so only translation by two is possible, first at
$T/2$. On the parity-twisted $C_{2p}$ with $p$ odd only primitive
frequencies are active, and a translation fixing the frequencies $2$ and
$p$ has displacement divisible by $2p$, so it is the identity.
\end{proof}

\subsection{Optimal signings of square cycles}\label{gra:sec:cycles}
\sources{grh_transfers_cycle_signings_and_lubin_tate; grh_transfer_advances_2026-11-26}

What is the smallest possible spectral radius of the adjacency matrix of
the square of a cycle, when each edge can be signed independently? For
$n\ge5$ let $C_n(1,2)$ have vertex set $\Z/n\Z$, with an undirected edge
between vertices at cyclic distance one or two. Write $a_i\in\{\pm1\}$ for
the sign of $\{i,i+1\}$ and $b_i\in\{\pm1\}$ for that of $\{i,i+2\}$. The
signed adjacency matrix $A$ is real and symmetric, and
$\rho(A)=\max_{\lambda\in\operatorname{spec}(A)}\abs\lambda=\norm A$. Switching
signs at vertices conjugates $A$ by a diagonal matrix with entries $\pm1$
and preserves the spectrum. The \emph{triangle signs} and the
\emph{holonomy} (the Hamilton-cycle sign product)
\[
 \tau_i=a_ia_{i+1}b_i,\qquad\alpha=\prod_{i\bmod n}a_i
\]
are invariant under switching and determine a switching class on labelled
vertices: choose vertex signs successively to set $a_0,\dots,a_{n-2}$
equal to $+1$; then $a_{n-1}=\alpha$ and $b_i=\tau_ia_ia_{i+1}$. There are
$2^{n+1}$ such classes. Counts of minimising switching classes below are
not counts of individual edge assignments.

Suvagiya's Conjecture~28 proposed the constant optimum
$r_*=F(2)$ for $n=8m\ge32$, in the notation of \cref{gra:thm:cycle8m}
\cite{Suvagiya2026Parity}. A counterexample at $n=32$ attained a value
below $2.79$ and left the exact optimum open
\cite{Trureturing2026Refutation}. The results below give the exact
optimum for every $n=8m\ge32$ and for every even $n$ with $8\le n\le50$.

\paragraph{Path compressions.}
For a word $w=(t_0,\dots,t_{k-1})$ of signs, let $H(w)$ be the signed
square of a path on $k+2$ vertices, with all length-one signs positive and
length-two signs $t_i$. An interval of $k+2$ consecutive vertices of
$C_n(1,2)$ induces such a path when $n\ge k+4$: its endpoints are then at
cyclic distance at least three, so there is no extra seam edge, and a
switch on the interval yields $H(w)$ for the corresponding consecutive
triangle signs. The compression of a real symmetric matrix to a coordinate
subspace has norm at most the original norm, by the Rayleigh quotient.
Hence
\begin{equation}\label{gra:eq:compress}
 \rho(A)<\gracut\ \Longrightarrow\ \norm{H(w)}<\gracut
 \ \text{for every window $w$ of the triangle word of length at most $n-4$.}
\end{equation}

A word passes the test $\norm{H(w)}<\gracut$ precisely when both
$\gracut I-H(w)$ and $\gracut I+H(w)$ are positive definite; then each
leading principal submatrix is positive definite, so every prefix passes.
Starting with the empty word and appending each of the two signs to each
surviving word is therefore exhaustive. After eliminating all but the last
two vertices of $\gracut I-H(w)$, write the remaining Schur block as
$\bigl(\begin{smallmatrix}a&b\\b&c\end{smallmatrix}\bigr)$. Appending a
sign $t$ adjoins a boundary column with entries $-t,-1$, and eliminating
the vertex with pivot $a$ gives
\begin{equation}\label{gra:eq:schur}
 (a,b,c)\longmapsto\Bigl(c-\frac{b^2}{a},\ -1+\frac{bt}{a},\
 \gracut-\frac1a\Bigr).
\end{equation}
The extension is kept if $a'>0$ and $a'c'>(b')^2$ for both shifted
matrices. Alternating vertex signs transform $\gracut I+H(w)$ into
$\gracut I-H(-w)$, so the second recurrence uses $-t$. Both start at
$(\gracut,-1,\gracut)$. For rational $\gracut$ this is an exact rational
computation with no eigenvalue tolerance.

Put
\[
 \graPi=(+,+,-,+,-,-,+,-).
\]

\begin{lemma}[Finite obstruction]\label[lemma]{gra:lem:windows}
Let $\gracut=1397/500$. Exactly eight words $w$ of length $27$ have
$\norm{H(w)}<\gracut$: the eight length-$27$ windows of $\graPi$ repeated
periodically. At length $19$ the only additional words are the two
alternating windows. The numbers of surviving words of each length are
\begin{center}\small\upshape
\begin{tabular}{r|rrrrrrrrrrrrrr}
\toprule
Length &0&1&2&3&4&5&6&7&8&9&10&11&12&13\\
Survivors &1&2&4&6&10&16&24&28&36&42&46&50&42&42\\
\midrule
Length &14&15&16&17&18&19&20&21&22&23&24&25&26&27\\
Survivors &44&38&32&22&18&10&10&10&10&10&10&10&10&8\\
\bottomrule
\end{tabular}
\end{center}
\end{lemma}

\begin{proof}
This is the evaluation of the recurrence \eqref{gra:eq:schur}; a prefix
failing either strict inequality can have no admissible descendant. Counts
alone do not identify the surviving words, and the final set comparison
\[
 \{w:\abs w=27,\ \norm{H(w)}<\gracut\}
 =\bigl\{(\graPi_{j+i\bmod8})_{0\le i<27}:0\le j<8\bigr\}
\]
is made explicitly, as is the comparison of the length-$19$ words with the
indicated union. The certificates are
\grafiles{\artifact{grh_transfers_cycle_signings_and_lubin_tate/graph_certificate.txt}\\
\artifact{grh_transfers_cycle_signings_and_lubin_tate/permutation_certificate.txt}}
The first also contains a second calculation, which builds the integer
matrices $1397I\pm500H(w)$ with the actual signs; for every surviving
parent, positivity of the two new determinants is equivalent to both
leading-minor tests, since all earlier leading minors are already
positive, and the determinants are recomputed by fraction-free integer
elimination. The second is an independent implementation that sums signed
permutation products and obtains the same eight words.
\end{proof}

\begin{proposition}[Rigidity below the cutoff; \stU]\label[proposition]{gra:prop:rigidity}
Let $\rho(A)<1397/500$. If $n\ge31$, then $8\mid n$ and the triangle signs
form a rotation of repeated $\graPi$. If $n\ge16$, the triangle word is
either a rotation of repeated $\graPi$, and then $8\mid n$, or alternating.
\end{proposition}

\begin{proof}
Put $\gracut=1397/500$. By \eqref{gra:eq:compress} every length-$27$ window
must be in the finite list if $n\ge31$. Its first $26$ signs determine its
phase, since the eight rotations already have distinct first eight signs.
Overlap forces the next letter at every step. The primitive period is
eight, so closing the word around the cycle requires $8\mid n$.

For $n\ge16$, regard each of the $42$ surviving length-$12$ words of
\cref{gra:lem:windows} as an arrow from its $11$-letter prefix to its
$11$-letter suffix. Delete any arrow whose initial vertex has no incoming
arrow or whose final vertex has no outgoing arrow; no arrow on a closed
walk can be removed. The counts after successive rounds are $42,20,10,10$,
and the ten remaining arrows form the period-eight cycle of $\graPi$ and
the alternating cycle (second certificate of \cref{gra:lem:windows}). Each
length-$12$ word occupies $14$ vertices, so by \eqref{gra:eq:compress}
every cyclic length-$12$ window of the triangle word survives when
$n\ge16$, and the cyclic word is a closed walk on the ten remaining arrows.
\end{proof}

At $n=32$ there are $2^{33}$ labelled switching classes before the test.
The proposition leaves sixteen possibilities below the threshold, one
rotation for each choice of $\alpha$.

\paragraph{The two boundary conditions.}
Normalise the length-one edges on the covering line. For repeated $\graPi$
the adjacency acts by
\[
 (Af)_i=f_{i-1}+f_{i+1}+\graPi_{i-2}f_{i-2}+\graPi_if_{i+2},
\]
indices of $\graPi$ being read modulo eight. The condition after one
circuit of $C_{8m}(1,2)$ is $f_{i+8m}=\alpha f_i$. Diagonalise translation
by eight: on the block indexed by $z$ with $z^m=\alpha$, $f_{i+8}=zf_i$.
The $m$ blocks are mutually orthogonal because
$\sum_{j=0}^{m-1}(z/z')^j=0$ for $z\ne z'$, and their dimensions add to
$8m$. Suvagiya's equation~(3) gives the characteristic polynomial of the
eight-vertex block \cite{Suvagiya2026Parity}, and the same symbol applies
to either holonomy; it is also obtained by building the symbol from the
signed edges with $z$ an indeterminate and expanding its determinant (this
check is part of the first certificate of \cref{gra:lem:windows}).
With $u=z+z^{-1}$ and $h=x^2-4$ the polynomial is
\[
 h^4-(16+2u)h^2+u^2+19u+38 .
\]
The two values of $h^2$ are $8+u\pm\sqrt{26-3u}$. Taking the greatest
possible value of $x^2=4+h$, the block norm is
\[
 F(u)=\sqrt{4+\sqrt{8+u+\sqrt{26-3u}}} .
\]
The function $F$ increases strictly on $[-2,2]$, since
\[
 \frac{\dd}{\dd u}\bigl(8+u+\sqrt{26-3u}\bigr)=1-\frac3{2\sqrt{26-3u}}>0 .
\]
In the periodic case $z=1$ attains $u=2$. In the antiperiodic case the
largest allowed value is $2\cos(\pi/m)$.

\begin{theorem}[Unrestricted minimum for $C_{8m}(1,2)$ and all minimisers; \stU]\label{gra:thm:cycle8m}
For every $m\ge4$ the minimum of $\rho(A)$ over all signings of
$C_{8m}(1,2)$ is
\[
 F(2\cos(\pi/m)).
\]
Precisely eight labelled switching classes attain it, with triangle word
a rotation of $\graPi^m$ and $\alpha=-1$. The next distinct value is
$r_*=F(2)$, attained by eight classes with the same rotations and
$\alpha=+1$. All remaining classes have $\rho(A)\ge1397/500$.
\end{theorem}

\begin{proof}
The two candidate radii are strictly ordered by monotonicity. Both are
below $1397/500$: $\sqrt{20}<4.473$, hence
$F(2)^2=4+\sqrt{10+\sqrt{20}}<7.805<(1397/500)^2$.
\Cref{gra:prop:rigidity} excludes every other word from this interval.
Completeness of the switching invariants gives the exact counts.
\end{proof}

For $n=32$, $u=\sqrt2$ at the winning phase and the minimum is
$2.78426979930755\ldots$; the optimiser is precisely the switching orbit,
up to rotation, of the counterexample of
\cite{Trureturing2026Refutation}, whose optimality now follows for
arbitrary choices on all $64$ edges. As $m$ increases, the minimum
increases strictly to $r_*$, with
$r_*-F(2\cos(\pi/m))=\pi^2F'(2)m^{-2}+O(m^{-4})$. Thus the conjectured
constant $r_*$ is the second distinct spectral radius.

\paragraph{The alternating family.}
For the alternating triangle word and even $n$, the Fourier blocks at
negative holonomy are
\[
 \begin{pmatrix}2\cos2\theta&2\cos\theta\\2\cos\theta&-2\cos2\theta
 \end{pmatrix},\qquad\theta=(2j+1)\pi/n,
\]
with squared norm $4(\cos^2\theta+\cos^22\theta)$. Putting
$x=\cos^2\theta$ makes this expression divided by four equal to
$4x^2-3x+1$, a convex polynomial. For even $n\ge8$ its maximum on the
allowed values is at $x_{\max}=\cos^2(\pi/n)$: at the opposite extreme the
allowed value is $0$ or $\sin^2(\pi/n)$. If it is $0$, $x_{\max}>3/4$
implies $4x_{\max}^2-3x_{\max}+1>1$. If it is $1-x_{\max}$, subtracting
the two polynomial values gives $2x_{\max}-1>0$. Hence the alternating
signing with $\alpha=-1$ has spectral radius
\begin{equation}\label{gra:eq:alternating}
 G(n)=2\sqrt{\cos^2(\pi/n)+\cos^2(2\pi/n)},
\end{equation}
while the positive-holonomy radius is $2\sqrt2$. (These are the values of
\cite[Proposition~25]{Suvagiya2026Parity}.) The function $G$ is
increasing.

\begin{theorem}[Even lengths at which alternation is optimal; \stU]\label{gra:thm:alternating}
For every even $n$ with $8\le n\le30$ and for
$n\in\{34,36,38,42,44,46\}$, the unrestricted minimum of $\rho(A)$ over
all signings of $C_n(1,2)$ is $G(n)$. For
$n\in\{34,36,38,42,44,46\}$ exactly two labelled switching classes attain
it: alternating triangle signs with $\alpha=-1$.
\end{theorem}

In particular $n=32$ is the first even length at which alternation loses.
The values at $n=26,28,30$, for which \cite[Remark~29]{Suvagiya2026Parity}
records heuristic searches, are $2.777192$, $2.784204$ and $2.789872$ to
six decimals.

\begin{proof}
\emph{Lengths $16\le n\le30$.} Let $\gracut=1397/500$. By
\cref{gra:prop:rigidity}, a signing with $\rho(A)<\gracut$ has
period-eight or alternating triangle word. The increasing function $G(n)$ is below $\gracut$ for $n\le30$: put
$v=3.14/30$ and $C(t)=1-t^2/2+t^4/24$; then
$4(C(v)^2+C(2v)^2)<\gracut^2$ by integer arithmetic, and $C(v),C(2v)$
bound the cosines at $n=30$ from above. Since $2\sqrt2>\gracut$, the
alternating word contributes exactly $G(n)$. Period eight is compatible
only with $n=16,24$ in this range. At $16$,
$G(16)^2=4+\sqrt{2+\sqrt2}+\sqrt2<7.3$, while $F(0)^2>7.6$. At $24$,
$G(24)^2=4+(\sqrt6+\sqrt2)/2+\sqrt3<7.67$, while $F(1)^2>7.7$.

\emph{Lengths $8,10,12,14$.} These have direct finite certificates: with
cutoffs $2327/1000$, $1249/500$, $519/200$ and $26551/10000$ respectively,
only alternating triangle signs with negative holonomy survive:
\grafiles{\artifact{grh_transfers_cycle_signings_and_lubin_tate/small_cycle_certificate.txt}}

\emph{The six further lengths.} For each $n$ choose the rational
threshold $\gracut$ in the table below. The recursion
\eqref{gra:eq:schur} lists every word of length $n-4$ whose path norm is
below $\gracut$. Complete each such word in every one of the $2^4$ ways,
and retain a cycle word only when every cyclic window of length $n-4$ is
on the passing list. In each case the sole surviving rotation class is
alternating.
\begin{center}\small\upshape
\begin{tabular}{rrrrr}
\toprule
$n$&$\gracut$&window length&passing path words&cyclic rotation classes\\
\midrule
34&$281/100$&30&98&1\\
36&$281/100$&32&108&1\\
38&$281/100$&34&92&1\\
42&$281/100$&38&70&1\\
44&$703/250$&40&394&1\\
46&$703/250$&42&422&1\\
\bottomrule
\end{tabular}
\end{center}
This necessary condition is complete as an exclusion argument: a window
occupies $n-2$ consecutive vertices, and the two omitted consecutive
vertices leave its endpoints at cyclic distance three, so
\eqref{gra:eq:compress} applies. After reduction to rotation
representatives, both values of $\alpha$ are tested by exact integer
elimination on $(\operatorname{den}\gracut)(\gracut I\pm A)$: Bareiss
elimination computes the leading principal minors without row swaps, and
their positivity is Sylvester's criterion. Only $\alpha=-1$ passes, in
accordance with $2\sqrt2>\gracut$. Every prefix count is recorded in
\grafiles{\artifact{grh_transfer_advances_2026-11-26/checks/graphs/even_extension_certificate.json}}
Finally $G(n)<\gracut$ is certified
using a rational lower bound for $\pi$ from Machin's identity and the
alternating arctangent series, followed by
$\cos x\le1-x^2/2+x^4/24$. Hence the candidate exists below the threshold
that excludes all competitors, and the two phases of alternation give
exactly two labelled switching classes.
\end{proof}

\begin{theorem}[The fifty-vertex optimum; \stU]\label{gra:thm:fifty}
Let $\rho_{50}$ be the minimum spectral radius over all independent edge
signings of $C_{50}(1,2)$, and put
\[
 \Omega=(-\graPi)^6(-,+).
\]
Precisely the switching classes with $\alpha=+1$ and triangle word a
cyclic rotation of $\Omega$ or $-\Omega$ minimise the spectral radius.
These are $100$ distinct labelled switching classes. Moreover
\[
 2.811568906156<\rho_{50}<2.811568906157,
\]
and $\rho_{50}=\sqrt{y_{\max}}$, where $y_{\max}$ is the largest real root
of the degree-$25$ polynomial $Q$ of \cref{gra:tab:q50}. Every other
switching class has spectral radius at least $7029/2500=2.8116$.
\end{theorem}

\begin{proof}
Set $\gracut=7029/2500$. By the recurrence \eqref{gra:eq:schur}, in
rational arithmetic, there are $324$ passing words at triangle length
$46$. Append all four remaining signs to each of them and retain a
length-$50$ word only if all its cyclic length-$46$ windows pass. There
are exactly two resulting rotation classes, represented by $\Omega$ and
$-\Omega$. The local path classification is checked by two independent
algorithms, the Schur recurrence and the summation of signed permutation
products using only integer arithmetic and the band width of the actual
path matrices; they agree on every prefix count and every surviving word:
\grafiles{\artifact{grh_transfer_advances_2026-11-26/checks/graphs/n50/certificate_n50.json}\\
\artifact{grh_transfer_advances_2026-11-26/checks/graphs/n50/verification_manifest.json}}

Any $48$ consecutive vertices of $C_{50}(1,2)$ induce such a path: the two
endpoints have cyclic distance three, so no extra edge joins them. By
\eqref{gra:eq:compress} every signing with radius $<\gracut$ must have one
of these two triangle words, up to rotation. There remain just two
holonomies for each word. Exact fraction-free elimination of the full
matrices $7029I\pm2500A$ shows that $\alpha=+1$ passes both Sylvester
tests, whereas $\alpha=-1$ fails: the first $49$ leading minors are
positive in each failed test, and the final determinant is negative. Thus
no additional signing can have radius $<\gracut$.

Negating every edge changes every triangle sign to its negative and,
because $50$ is even, preserves $\alpha$. It takes $A$ to $-A$, so the two
surviving rotation classes have equal spectral radius, and both are
global minimisers. Each word has primitive period $50$ and their rotation
orbits are disjoint, giving exactly $100$ labelled switching classes.

The characteristic polynomial of either minimising representative is
$\det(xI-A)=Q(x^2)$, computed from the integer traces $\tr(A^j)$ by
Newton's identities. For $r_-=2811568906156/10^{12}$ and
$r_+=2811568906157/10^{12}$, exact leading-minor sign counts show that
$r_-I-A$ and $r_-I+A$ each have exactly one negative eigenvalue, while
both shifts at $r_+$ have none:
\grafiles{\artifact{grh_transfer_advances_2026-11-26/checks/graphs/n50/root_isolation.json}}
This proves the strict interval and shows that the interval contains
exactly one positive eigenvalue. No irreducibility assertion about $Q$ is
needed.
\end{proof}

\begin{table}[htbp]
\centering\small
\begin{tabular}{rr@{\qquad}rr}
\toprule
$j$&$q_j$&$j$&$q_j$\\
\midrule
25&1&12&$-42556190827574$\\
24&$-100$&11&95424283730427\\
23&4702&10&$-173149057931144$\\
22&$-138186$&9&251857479144425\\
21&2846105&8&$-290088695299758$\\
20&$-43660212$&7&260508422825773\\
19&517548579&6&$-178938476667148$\\
18&$-4854567112$&5&91816438267915\\
17&36601189842&4&$-34160308422738$\\
16&$-224095627758$&3&8848896828974\\
15&1121189638159&2&$-1499411878484$\\
14&$-4598060887914$&1&148268409185\\
13&15463123646497&0&$-6447640168$\\
\bottomrule
\end{tabular}
\caption{Coefficients of $Q(y)=\sum_{j=0}^{25}q_jy^j$, where
$\det(xI-A)=Q(x^2)$ for a minimising signing of $C_{50}(1,2)$.}
\label{gra:tab:q50}
\end{table}

\begin{remark}\label[remark]{gra:rem:fifty}
The best value in the alternating family at $n=50$ is $G(50)>2.814$: with
$\pi<22/7$, $\cos u\ge1-u^2/2$ and $v=11/175$, rational arithmetic gives
$4\bigl((1-v^2/2)^2+(1-2v^2)^2\bigr)>(2814/1000)^2$. Thus the improvement
over alternation at fifty vertices is strict. The winning word is six
copies of an eight-letter block with an additional two-letter segment, and
its holonomy is positive, whereas the minimisers at $n=8m\ge32$ have exact
period eight and negative holonomy. One cannot predict the optimiser by
carrying over the best boundary condition from a different cycle length.
\end{remark}

\begin{remark}[Settled cases]\label[remark]{gra:cor:settled}
By \cref{gra:thm:alternating,gra:thm:cycle8m,gra:thm:fifty}, the exact
minimum of $\rho(A)$ over all signings of $C_n(1,2)$ is known for every even
$n$ with $8\le n\le50$ and for every multiple of eight $n\ge32$; the
following table only collects these theorems:
\begin{center}\small\upshape
\begin{tabular}{lll}
\toprule
$n$ & minimum & minimising switching classes\\
\midrule
even, $8\le n\le30$ & $G(n)$ & alternating, $\alpha=-1$\\
$8m$, $m\ge4$ & $F(2\cos(\pi/m))$ & eight: rotations of $\graPi^m$, $\alpha=-1$\\
$34,36,38,42,44,46$ & $G(n)$ & two: alternating, $\alpha=-1$\\
$50$ & $\rho_{50}=2.811568906156\ldots$ & $100$: rotations of $\pm\Omega$, $\alpha=+1$\\
\bottomrule
\end{tabular}
\end{center}
\end{remark}

\begin{proposition}[A local gap and the infinite line; \stU]\label[proposition]{gra:prop:line}
At the cutoff $281/100$ there are $72$ surviving words of length $43$, and
their central nine-letter subwords (indices $17$ through $25$ from zero)
are exactly the eight windows of $\graPi$. Consequently:
\begin{enumerate}[label=\textup{(\alph*)}]
\item for $n\ge47$, any signing of $C_n(1,2)$ whose triangle word differs
from the period-eight family has $\rho(A)\ge2.81$;
\item the same holds for signings of the graph on $\Z$ joining integers at
distance one or two, and on this graph the minimum of $\norm A$ is
$r_*=F(2)$, attained uniquely up to switching and translation by the
period-eight triangle word. The minimum is isolated.
\end{enumerate}
\end{proposition}

\begin{proof}
The enumeration is recorded in the second certificate of
\cref{gra:lem:windows}. Part (a) follows from \eqref{gra:eq:compress}, since every position of
the cycle is central in a length-$43$ window. On the line every finite
compression is an operator-norm lower bound, and every position is
central in a length-$43$ window, which gives the gap. All length-one
signs on the line can be switched to $+1$, and Fourier transformation by
period-eight cells gives the norm of the surviving operator as
$\max_{-2\le u\le2}F(u)=r_*$.
\end{proof}

\begin{corollary}[Poles of the signed Ihara function; \stU]\label[corollary]{gra:cor:ihara}
For a signed four-regular graph, the Bass determinant formula for the
signed Ihara zeta function gives the nontrivial factor of its reciprocal
as
\[
 Q_A(v)=\det(I-vA+3v^2I)=\prod_{\lambda\in\operatorname{spec}(A)}
 (1-\lambda v+3v^2),
\]
the universal factor $(1-v^2)^{\abs E-\abs V}$ being omitted. Among all
signings of $C_{8m}(1,2)$, $m\ge4$, the largest possible minimum angular
distance of the roots of $Q_A$ from the real axis is
\[
 \arccos\Bigl(\frac{F(2\cos(\pi/m))}{2\sqrt3}\Bigr),
\]
and \cref{gra:thm:cycle8m} characterises every maximising signing. At
$n=32$ the angle is about $36.51^\circ$.
\end{corollary}

\begin{proof}
The displayed form also follows by taking the ordinary Bass identities for
the double cover and for the base graph and dividing them: adjacency on
the cover splits into the base and signed matrices
\cite{Suvagiya2026Parity,Bass1992Ihara}. When
$\abs\lambda\le2\sqrt3$ the two roots of the factor
$1-\lambda v+3v^2$ are $3^{-1/2}e^{\pm i\theta}$, where
$2\sqrt3\cos\theta=\lambda$; for $\abs\lambda>2\sqrt3$ there is a root on
the real axis. The minimum angular distance is therefore a decreasing
function of $\rho(A)$. The signed Ihara function itself has poles at
these points.
\end{proof}

\subsection{A bipartite family with all zeros on the canonical line}\label{gra:sec:ehrhart}
\sources{grh_method_transfers}

For a graph $G$ on $v$ vertices let $P_G$ be the convex hull of
$\pm(e_i-e_j)$ over its edges, in the lattice
$\Z^v\cap\{\sum x_i=0\}$. Write
$E_{m,n}(k)=\abs{kP_{K_{m,n}}\cap\Z^{m+n}}$ for the Ehrhart polynomial of
the complete bipartite graph $K_{m,n}$; its degree is $d=m+n-1$.
Higashitani, Kummer and Micha\l{}ek asked, for complete multipartite
graphs, whether all zeros lie on the line $\Re z=-1/2$ and whether
decreasing a largest part gives interlacing
\cite[Conjecture~4.10]{HigashitaniKummerMichalek2016Conjecture}. A
counterexample with six parts of size two has been reported
\cite{DavisCounterexampleEhrhart}; it leaves the complete bipartite question intact.
K\"olbl gave a conditional result for $K_{4,n}$
\cite[Theorem~3.3]{Kolbl2024Ehrhart}, and the line property through
smaller part five was proved using positive Jacobi matrices
\cite{Carptopus2026Ehrhart}. Here the transfer from the source proof is
methodological: after rotation, exclusion of zeros from one half-plane
suffices by symmetry, and the needed positivity comes from the known
Ehrhart expansion and finite certificates.

\begin{theorem}[Finite-width theorem; \stU]\label{gra:thm:ehrhart}
The following hold for the Ehrhart polynomials $E_{m,n}$.
\begin{enumerate}[label=\textup{(\alph*)}]
\item For $2\le m\le32$ and every integer $n\ge m$, all zeros of $E_{m,n}$
are simple and lie on $\Re z=-1/2$, and the ordinates of the zeros
strictly alternate with those of $E_{m,n-1}$.
\item For $2\le m\le32$ and $n\ge m$, the ordinates of the zeros of
$E_{m,n}$ strictly interlace those of $E_{m+1,n}$. Consequently all zeros
of $E_{33,n}$, $n\ge33$, are simple and lie on $\Re z=-1/2$.
\end{enumerate}
\end{theorem}

Width $32$ is the last width for which a positive Bezoutian is certified
below; consecutive-$n$ interlacing within width $33$ remains outside the
certificate. In particular, the interlacing of $E_{3,n}$ and $E_{4,n}$
holds for $n\ge3$; the cases $n=1,2$ of this particular interlacing follow
from the cross-polytope recurrence and from the larger-part interlacing at
width two, respectively. Width one follows at once from the cross-polytope
recurrence.

\paragraph{Reduction to two short polynomials.}
Put $r=m-1$, so that $d=n+r$, and let $C_k$ denote the Ehrhart polynomial
of a $k$-dimensional cross-polytope; its Ehrhart series is
$(1+t)^k/(1-t)^{k+1}$. Set $Q_k(y)=k!\,i^{-k}C_k((iy-1)/2)$. The
cross-polytope recurrence gives
\[
 Q_0=1,\qquad Q_1=y,\qquad Q_{k+1}=yQ_k-k^2Q_{k-1}.
\]
In particular $Q_k$ is the characteristic polynomial of a real symmetric
tridiagonal matrix with consecutive off-diagonal entries $1,2,\dots,k-1$.
The formula of Higashitani, Jochemko and Micha\l{}ek for the
$h^*$-polynomial
\cite[Theorem~4.1]{HigashitaniJochemkoMichalek2018Hstar} has
coefficients $\gamma_j=\binom{2j}{j}\binom rj\binom{n-1}j$ in the basis
$t^j(1+t)^{d-2j}$. Use $4t=(1+t)^2-(1-t)^2$ in the Ehrhart series.
Binomial expansion expresses $4^rE_{m,n}$ as
$\sum_{h=0}^r(-1)^hc_hC_{d-2h}$, where
$c_h=\sum_{j=h}^r4^{r-j}\gamma_j\binom jh$. Thus the rotated polynomial
has the exact expression
\begin{equation}\label{gra:eq:rotated}
 R_n(y)=4^rd!\,i^{-d}E_{m,n}((iy-1)/2)
 =\sum_{h=0}^rc_h(d)_{2h}Q_{d-2h}(y)
 =A_n(y)Q_n(y)-B_n(y)Q_{n-1}(y),
\end{equation}
where $(d)_k=d(d-1)\cdots(d-k+1)$, and where $A_n,B_n$ are defined by
expanding each $Q_{d-2h}$ from $Q_n,Q_{n-1}$ by forward or backward
recurrence. Their $y$-degrees are at most $r$ and $r-1$; $A_n$ has degree
$r$ with positive leading coefficient $c_0$; and their parities are those
of $r$ and $r-1$ respectively. The construction also works at $n=m-1$,
including a possible $Q_0$.

\paragraph{The positive form and the interlacing determinant.}
Let $A=A_n$, $B=B_n$, and let $\mathsf H$ be the real symmetric matrix of
coefficients in $u^iv^j$ ($0\le i,j<r$) of the Bezoutian
\[
 \frac{A(u)B(v)-B(u)A(v)}{u-v}.
\]

\begin{lemma}\label[lemma]{gra:lem:bezout}
Suppose $\mathsf H$ is positive definite. Then all zeros of $R_n$ are real
and simple. If moreover the even polynomial
\[
 T_n(y)=n^2A_nA_{n+1}-yB_nA_{n+1}+B_nB_{n+1}
\]
is positive on the real axis, then the zeros of $R_{n+1}$ are real and
simple and strictly interlace those of $R_n$.
\end{lemma}

\begin{proof}
Substitution of $v=\bar u$ shows that $A,B$ have no common zero and that
$\Im(B(u)/A(u))<0$ for $\Im u>0$ where the quotient is defined. In fact
$A(u)$ cannot vanish there, since the diagonal Hermitian form at
$(1,u,\dots,u^{r-1})$ would vanish. On the upper half-plane
$Q_n/Q_{n-1}$ has positive imaginary part, by induction using
$Q_{k+1}/Q_k=y-k^2/(Q_k/Q_{k-1})$. The two ratios cannot coincide, so
$R_n$ has only real zeros. On the real axis the diagonal limit gives
$W=A'B-AB'>0$. Also $W_n=Q_n'Q_{n-1}-Q_nQ_{n-1}'>0$, since $W_1=1$ and
$W_n=Q_{n-1}^2+(n-1)^2W_{n-1}$ for $n\ge2$.

Set $C=yA_{n+1}-B_{n+1}$ and $D=n^2A_{n+1}$, so that
$R_{n+1}=CQ_n-DQ_{n-1}$ and $T_n=AD-BC$. At every real root $t$ of $R_n$
there is a nonzero $\lambda$ with
$(Q_n,Q_{n-1})(t)=\lambda(B,A)(t)$; both vectors are nonzero by the
recurrence and by positive definiteness. Differentiation gives
$R_n'(t)=(\lambda^2W(t)+W_n(t))/\lambda$, which proves simplicity.
Moreover
\[
 \frac{R_{n+1}(t)}{R_n'(t)}=-\frac{\lambda^2T_n(t)}{\lambda^2W(t)+W_n(t)} .
\]
When $T_n$ is positive on the real axis, signs force a root of $R_{n+1}$
in each interior gap and one in each exterior gap, since both leading
coefficients are positive. The degree difference is one, so this accounts
for all its roots, with multiplicity one.
\end{proof}

\begin{proof}[Proof of \cref{gra:thm:ehrhart}\textup{(a)}]
By \eqref{gra:eq:rotated} and \cref{gra:lem:bezout} it suffices to prove,
for each $2\le m\le32$, that $\mathsf H$ is positive definite for every
$n\ge m-1$ and that $T_n$ is positive on the real axis for every
$n\ge m-1$. Only short polynomials involving $r$ occur in the two
positivity tests, and finite exact calculations prove the universal
quantifier over $n$, as follows. The integer calculations for widths $2$
through $32$ are recorded in
\artifact{grh_method_transfers/checks/ehrhart_output.txt} and in
\artifact{grh_method_transfers/checks/interlacing_output.txt}.

\emph{Degree bounds.} In performing the backward recurrence, multiply
both coordinates by $S(n)=\prod_{j=1}^{r-1}(n-j)^2$. Backward steps divide
by distinct squared linear factors of $S$. In the forward direction at
most $r$ steps increase the degree in $n$ by at most $2r$. In the backward
direction the coordinates have denominators dividing $S$ and, after
multiplication by $S$, degrees at most $2r$ (each step is
$[yq_{k+1}-q_{k+2}]/(n+k+1)^2$, where indices are relative to $n$). Thus
each cleared coordinate in the mixture has degree at most $7r$ in $n$:
$2r$ allows for $S$, $2r$ for the forward recurrence, $r$ for $c_h$, and
$2r$ for $(d)_{2h}$. At $7r+1$ consecutive integers beginning at $m$, one
verifies by taking forward differences that the coordinates divided by $S$
agree with polynomials of degrees at most $r$ for $A_n$ and $r+1$ for
$B_n$. Multiplying these interpolants by $S$ and applying the polynomial
identity principle proves the degree bounds identically.

\emph{The Bezoutian.} The Bezoutian splits into even and odd blocks by
parity. Each entry has degree at most $2r+1$ in $n$. For the leading minor
of order $j$ of either block, its exact integer values at
$n=m,m+1,\dots,m+(2r+1)j$ are computed by fraction-free elimination. All
forward differences at $m$, through its actual degree, are positive. For
every integer $s\ge0$, Newton's identity
$f(m+s)=\sum_h\binom sh\Delta^hf(m)$ now proves strict positivity, and
Sylvester's criterion applies to both blocks. At $n=m-1$ the two blocks
are tested directly.

\emph{The interlacing determinant.} Write $T_n(y)=U(n,y^2)$; its degree
in $n$ is at most $2r+2$. The differences $\Delta_n^hU(m,z)$ are obtained
from $2r+3$ rows. For each nonzero polynomial in $z$ among them one
verifies a positive constant term and equal numbers of Sturm sign changes
at $0$ and $+\infty$. Hence each is positive on $[0,\infty)$, and so is
$U(m+s,z)$ for integral $s\ge0$ by Newton's identity (the zeroth row is
positive). The same Sturm test at $n=m-1$ completes the square boundary
case.
\end{proof}

\begin{proof}[Proof of \cref{gra:thm:ehrhart}\textup{(b)}]
Keep the pair $Q_n,Q_{n-1}$ fixed, and use subscripts $(m,n)$ for the
short coordinates. In the residue formula of \cref{gra:lem:bezout} take
$C=A_{m+1,n}$ and $D=B_{m+1,n}$. The determinant
\[
 V(n,y^2)=A_{m,n}(y)B_{m+1,n}(y)-B_{m,n}(y)A_{m+1,n}(y)
\]
has even parity. Its degree in $n$ is at most $2r+2$, as the termwise
degree bounds for its factors are $(r,r+2)$ and $(r+1,r+1)$. The
coordinate bounds for width $m+1$ are tested by the same interpolation
identity, with upper endpoint including width $33$, and $V(m+s,z)>0$ is
certified for every integer $s\ge0$ and real $z\ge0$ by applying the Sturm
test to the Newton coefficient polynomials in $s$; the output is
\artifact{grh_method_transfers/checks/smaller_part_output.txt}. The
residue sign argument gives the claimed interlacing. This argument needs
the positive Bezoutian only in the lower width.
\end{proof}

\begin{example}\label[example]{gra:ex:k66}
For $K_{6,6}$, dividing both short coordinates by their common factor
$132$ gives
\[
 A(y)=693y^5-29984y^3+175507y,\qquad
 B(y)=19656y^4-442620y^2+881100 .
\]
Thus positivity of a five-by-five real form suffices for the
degree-eleven Ehrhart polynomial in this example. Its first four
lattice-point counts, at dilations $0,1,2,3$, are $1,73,1915,26547$.
\end{example}

\begin{proposition}[Membership in the dilated polytope; \stU]\label[proposition]{gra:prop:membership}
Split a lattice vector into its coordinates $(u,v)$ on the two sides of
$K_{m,n}$. It belongs to $kP_{K_{m,n}}$ if and only if
\[
 \textstyle\sum u_i+\sum v_j=0,\qquad
 \max(\norm u_1,\norm v_1)\le k .
\]
\end{proposition}

\begin{proof}
A weighted sum of signed edges is a matrix with row sums $u$, column sums
$-v$, and cost equal to the sum of the absolute entries. Both marginal
norms bound that cost from below. For the converse let $T$ be the larger
marginal norm and let $s=\sum u_i$. On each side choose nonnegative
positive and negative margins with total positive mass $(T+s)/2$ and total
negative mass $(T-s)/2$: starting from the ordinary positive and negative
parts, add the same padding to both where needed. Nonnegative matrices
with each pair of prescribed margins exist by successively matching
supplies and demands. The difference of the two matrices realises the
desired vector with cost at most $T$. Real edge weights suffice for
membership in the polytope.
\end{proof}

Counting integer vectors by their sum and $\ell^1$ norm on each side
recovers the entire polynomial from the required $d+1$ integer values for
$K_{2,2}$, $K_{3,5}$, $K_{6,6}$ and $K_{7,9}$, a total of $40$ comparisons
with \eqref{gra:eq:rotated}; see
\artifact{grh_method_transfers/checks/instances_output.txt}.

\subsection{Open questions}\label{grw:sec:open}
{\raggedright\sources{grh_method_transfers; grh_transfers_cycle_signings_and_lubin_tate; grh_transfer_advances_2026-11-26; grh_cayley_and_filters}}

\begin{enumerate}
\item \emph{Oriented Cayley walks.} For a prescribed conductor $f$ and
spacing $a\sqrt{-f}$, what is the least outdegree of a periodic oriented
Cayley graph that is connected on the whole host?
\Cref{gra:prop:cayleycoset} shows that this requirement can change the
answer of \cref{gra:thm:cayleyperiods}. For larger conductors, which
connection sets attain the degree bound on hosts of order greater than
$af$? \Cref{gra:cor:cayleyduality} constrains them by a duality, the
spectrum and the rank. Nonabelian groups would require control of
higher-dimensional irreducible representations.
\item \emph{Square cycles at lengths incompatible with period eight.}
What are the exact optima of \cref{gra:sec:cycles} at large lengths
incompatible with period eight? The first undetermined even length is
$n=52$. Is there a uniform
description at lengths $8m+2$, $8m+4$ and $8m+6$? The fifty-vertex word is
evidence for a defect-pattern search, not a proof that one defect always
wins, and the finite-window language can become much larger at a slightly
increased threshold.
\item \emph{A conceptual proof of period eight.} Can the $29$-vertex test
of \cref{gra:lem:windows} be replaced by a short human-checkable family
of inequalities for forbidden sign patterns? Such a proof would explain
the period eight and could extend to $C_n(1,k)$ with $k>2$.
\item \emph{All widths.} For complete bipartite edge polytopes, are the
two parity blocks of the Bezoutian of \cref{gra:lem:bezout} positive in
every width? A representation of their minors as sums with positive
coefficients could replace the width-by-width certificates and give a
route to the line property and to interlacing for unrestricted part
sizes.
\end{enumerate}
